% Artigo LNCS — Estudo Comparativo de Metaheurísticas (versão compacta ≤10 páginas)
% Compilar com: pdflatex artigo_v3.tex (duas vezes)
\documentclass[runningheads]{llncs}

\usepackage[T1]{fontenc}
\usepackage[utf8]{inputenc}
\usepackage[brazil]{babel}
\usepackage{graphicx}
\usepackage{amsmath}
\usepackage{amssymb}
\usepackage{booktabs}
\usepackage{multirow}
\usepackage{multicol}
\usepackage{array}
\usepackage{float}
\usepackage{algorithm}
\usepackage{algpseudocode}
\usepackage{hyperref}
\usepackage[table]{xcolor}
\usepackage{caption}
\usepackage{subcaption}

\newcolumntype{C}[1]{>{\centering\arraybackslash}p{#1}}
\algrenewcommand\algorithmicindent{0.8em}

\begin{document}

\title{Estudo Comparativo de Metaheurísticas para Otimização
de Fun\c{c}\~{o}es Benchmark: AG, DE, PSO e EPSO
em M\'{u}ltiplas Dimens\~{o}es}

\titlerunning{Estudo Comparativo de Metaheur\'{i}sticas}

\author{Davi\inst{1}}
\authorrunning{Davi}

\institute{Universidade Federal do Rio de Janeiro (UFRJ), Rio de Janeiro, Brasil\\
\email{aluno@ic.ufrj.br}}

\maketitle

%─────────────────────────────────────────────────────────────────────────────
\begin{abstract}
Este trabalho apresenta um estudo comparativo entre quatro algoritmos
metaheurísticos de otimização contínua: Algoritmo Genético~(AG),
Evolução Diferencial~(DE), Otimização por Enxame de Partículas~(PSO)
e Otimização Evolutiva por Enxame de Partículas~(EPSO). Os algoritmos
são avaliados em dez funções benchmark escaláveis — cinco unimodais e
cinco multimodais — nas dimensões 30, 50 e 100, com 51~execuções
independentes e orçamento fixo de 100.000~avaliações de função.
Os hiperparâmetros foram ajustados via Optuna/TPE\@. A comparação
estatística utiliza Teste~T de Welch com Wilcoxon Rank-Sum como
verificação complementar ($\alpha=5\%$). Os resultados evidenciam
que nenhum algoritmo domina em todos os cenários: PSO lidera em
funções separáveis e velocidade de convergência, enquanto AG
destaca-se em funções multimodais densas.

\keywords{Metaheurísticas \and AG \and DE \and PSO \and EPSO
\and Funções Benchmark \and Otimização Contínua}
\end{abstract}

%─────────────────────────────────────────────────────────────────────────────
\section{Introdução}
\label{sec:intro}

Problemas de otimização contínua são frequentemente não-convexos,
multimodais e de alta dimensão, tornando os métodos baseados em gradiente
inadequados~\cite{storn1997}. Metaheurísticas populacionais contornam
essas limitações por meio de busca estocástica, mas sua eficácia relativa
varia conforme a estrutura do problema — princípio formalizado pelo
teorema \emph{No Free Lunch}~\cite{wolpert1997}.

O AG~\cite{goldberg1989} representa o paradigma evolutivo clássico; a
DE~\cite{storn1997} introduziu mutação diferencial como mecanismo de busca
competitivo; o PSO~\cite{kennedy1995} modelou a inteligência coletiva de
enxames; o EPSO~\cite{miranda2002} estendeu o PSO com auto-adaptação
de pesos e perturbação do melhor global. Apesar da literatura extensa
sobre cada algoritmo isoladamente, comparações sistemáticas sob protocolo
uniforme com ajuste moderno de hiperparâmetros permanecem relevantes.
Este trabalho preenche essa lacuna avaliando os quatro algoritmos em
dez funções benchmark em múltiplas dimensões.

%─────────────────────────────────────────────────────────────────────────────
\section{Algoritmos}
\label{sec:algoritmos}

Os quatro algoritmos compartilham a mesma interface: população de
$N=100$ indivíduos, domínio $(l,u)^D$, orçamento maxFEs, e retornam
o melhor fitness encontrado junto ao histórico de convergência por FEs.

\subsection{Algoritmo Genético (AG)}

O AG evolui uma população por seleção por torneio ($k=3$), cruzamento
BLX-$\alpha$ — que gera descendentes no intervalo estendido
$[\min-\alpha\Delta,\,\max+\alpha\Delta]$, $\Delta=|p_1-p_2|$ —
e mutação \emph{creep} gaussiana de desvio~$\sigma$ por gene com
probabilidade~$p_m$. A avaliação é vetorizada por geração.

\begin{algorithm}[H]
\caption{Algoritmo Genético (AG)}\label{alg:ag}
\begin{multicols}{2}
\begin{algorithmic}[1]
\footnotesize
\Require $f,D,(l,u),N,\text{maxFEs},k,\alpha,p_m,\sigma$
\State $P\!\leftarrow\!\mathcal{U}(l,u)^{N\times D}$,
       $\text{fit}\!\leftarrow\!f(P)$,
       $\text{FEs}\!\leftarrow\!N$
\While{$\text{FEs}<\text{maxFEs}$}
  \For{$i=1$ \textbf{to} $N$}
    \State $p_1\!\leftarrow\!\text{Torneio}(P,\text{fit},k)$
    \State $p_2\!\leftarrow\!\text{Torneio}(P,\text{fit},k)$
    \State $o_i\!\leftarrow\!\text{BLX-}\alpha(p_1,p_2)$
    \State $o_i\!\leftarrow\!\text{Creep}(o_i,p_m,\sigma)$
  \EndFor
  \State $P\!\leftarrow\!O$;
         $\text{fit}\!\leftarrow\!f(P)$
  \State $\text{FEs}\!\leftarrow\!\text{FEs}+N$
\EndWhile
\State \Return $\arg\min_P \text{fit}$
\end{algorithmic}
\end{multicols}
\end{algorithm}

\subsection{Evolução Diferencial (DE)}

A DE/rand/1/bin gera o vetor mutante $v_i = x_{r_0}+F(x_{r_1}-x_{r_2})$
com $r_0,r_1,r_2$ distintos e aleatórios. O cruzamento binomial com
$j_\text{rand}$ garantido combina $v_i$ e $x_i$ gene a gene com
probabilidade CR\@. A seleção \emph{greedy} substitui $x_i$ somente se
$f(u_i)\leq f(x_i)$.

\begin{algorithm}[H]
\caption{Evolução Diferencial (DE/rand/1/bin)}\label{alg:de}
\begin{multicols}{2}
\begin{algorithmic}[1]
\footnotesize
\Require $f,D,(l,u),N,\text{maxFEs},F,\text{CR}$
\State $P\!\leftarrow\!\mathcal{U}(l,u)^{N\times D}$,
       $\text{fit}\!\leftarrow\!f(P)$,
       $\text{FEs}\!\leftarrow\!N$
\While{$\text{FEs}<\text{maxFEs}$}
  \For{$i=1$ \textbf{to} $N$}
    \State $v_i\!\leftarrow\!x_{r_0}\!+\!F(x_{r_1}\!-\!x_{r_2})$
    \State $v_i\!\leftarrow\!\text{clip}(v_i,l,u)$
    \State $j_r\!\leftarrow\!\mathcal{U}\{1,\ldots,D\}$
    \State $u_i\!\leftarrow\!\text{Crossover}(x_i,v_i,\text{CR},j_r)$
    \If{$f(u_i)\leq\text{fit}_i$}
      \State $x_i\!\leftarrow\!u_i$;
             $\text{fit}_i\!\leftarrow\!f(u_i)$
    \EndIf
    \State $\text{FEs}\!\leftarrow\!\text{FEs}+1$
  \EndFor
\EndWhile
\State \Return $\arg\min_P\text{fit}$
\end{algorithmic}
\end{multicols}
\end{algorithm}

\subsection{Otimização por Enxame de Partículas (PSO)}

Cada partícula atualiza velocidade segundo
$v_i \leftarrow w v_i + c_1 r_1(p_i^\text{best}-x_i)
+ c_2 r_2(g^\text{best}-x_i)$
com $r_1,r_2\sim\mathcal{U}(0,1)$ e $v_\text{max}=0{,}2(u-l)$.

\begin{algorithm}[H]
\caption{PSO}\label{alg:pso}
\begin{multicols}{2}
\begin{algorithmic}[1]
\footnotesize
\Require $f,D,(l,u),N,\text{maxFEs},w,c_1,c_2$
\State $X,V\!\leftarrow\!\mathcal{U}$;
       $\text{fit}\!\leftarrow\!f(X)$;
       $\text{FEs}\!\leftarrow\!N$
\State $\text{pbest}\!\leftarrow\!X$;
       $g^\text{best}\!\leftarrow\!\arg\min\text{fit}$
\While{$\text{FEs}<\text{maxFEs}$}
  \For{$i=1$ \textbf{to} $N$}
    \State $V_i\!\leftarrow\!wV_i
           +c_1r_1(p_i^\text{b}-X_i)
           +c_2r_2(g^\text{b}-X_i)$
    \State $V_i\!\leftarrow\!\text{clip}(V_i,\pm v_\text{max})$
    \State $X_i\!\leftarrow\!\text{clip}(X_i+V_i,l,u)$
    \State $\text{fit}_i\!\leftarrow\!f(X_i)$;
           $\text{FEs}\!+\!=\!1$
    \State Atualizar pbest, gbest
  \EndFor
\EndWhile
\State \Return $g^\text{best}$
\end{algorithmic}
\end{multicols}
\end{algorithm}

\subsection{EPSO}

O EPSO estende o PSO criando um clone de cada partícula com pesos
mutados $w_j^*=\text{clip}(w_j+\tau\mathcal{N}(0,1),0,1)$ e
perturbando o melhor global multiplicativamente
$g_j^*=g_j(1+w_\text{pert}\,\xi)$, $\xi\sim\mathcal{U}(0,1)$,
apenas nas dimensões da máscara de comunicação estocástica (prob.~$p_c$
por dimensão). A seleção \emph{greedy} mantém o melhor entre original
e clone, adotando os pesos do clone quando este vence.

\begin{algorithm}[H]
\caption{EPSO}\label{alg:epso}
\begin{multicols}{2}
\begin{algorithmic}[1]
\footnotesize
\Require $f,D,(l,u),N,\text{maxFEs},\tau,p_c$
\State Inicializar $X,V,W\in[0,1]^{N\times4}$,
       pbest, gbest; $\text{FEs}\!\leftarrow\!N$
\While{$\text{FEs}<\text{maxFEs}$}
  \For{$i=1$ \textbf{to} $N$}
    \State $W_i^*\!\leftarrow\!\text{clip}(W_i+\tau\mathcal{N}^4,0,1)$
    \State $g^*\!\leftarrow\!\text{PertGbest}(g^\text{b},W_i^*,p_c)$
    \State Mover clone com $W_i^*$ e original com $W_i$
    \State $\text{FEs}\!\leftarrow\!\text{FEs}+2$
    \If{$f(\text{clone})<f(\text{orig})$}
      \State $X_i\!\leftarrow\!\text{clone}$;
             $W_i\!\leftarrow\!W_i^*$
    \Else
      \State $X_i\!\leftarrow\!\text{orig}$
    \EndIf
    \State Atualizar pbest, gbest
  \EndFor
\EndWhile
\State \Return $g^\text{best}$
\end{algorithmic}
\end{multicols}
\end{algorithm}

%─────────────────────────────────────────────────────────────────────────────
\section{Fun\c{c}\~{o}es Benchmark}
\label{sec:benchmark}

As dez funções foram selecionadas para cobrir diferentes tipos de
dificuldade: separabilidade, multimodalidade, não-separabilidade,
condicionamento e deceptividade. Todas têm ótimo global $f^*=0$ e
são escaláveis para qualquer $D$. A Tabela~\ref{tab:funcoes} sumariza
as propriedades de cada função, e a Figura~\ref{fig:benchmark_grid}
mostra as superfícies de contorno em $D=2$.

\begin{table}[H]
\caption{Propriedades das dez funções benchmark.
U=unimodal, M=multimodal, S=separável, NS=não-separável.}
\label{tab:funcoes}
\centering
\scriptsize
\setlength{\tabcolsep}{3pt}
\begin{tabular}{llllp{5.2cm}}
\toprule
\textbf{Função} & \textbf{Tipo} & \textbf{Sep.} & \textbf{Domínio} & \textbf{Fórmula} \\
\midrule
Sphere       & U & S  & $[-100,100]^D$ & $\sum x_i^2$ \\
Rosenbrock   & U & NS & $[-5,10]^D$    & $\sum[100(x_{i+1}\!-\!x_i^2)^2\!+(x_i\!-\!1)^2]$ \\
Sum Squares  & U & S  & $[-10,10]^D$   & $\sum i\,x_i^2$ \\
Dixon-Price  & U & NS & $[-10,10]^D$   & $(x_1\!-\!1)^2\!+\!\sum i(2x_i^2\!-\!x_{i-1})^2$ \\
Zakharov     & U & NS & $[-5,10]^D$    & $\sum x_i^2 + (\sum 0{,}5i\,x_i)^2 + (\sum 0{,}5i\,x_i)^4$ \\
Rastrigin    & M & S  & $[-5{,}12,5{,}12]^D$ & $10D+\sum[x_i^2-10\cos(2\pi x_i)]$ \\
Ackley       & M & NS & $[-32{,}77,32{,}77]^D$ & $-20e^{-0{,}2\sqrt{\frac{\sum x_i^2}{D}}}-e^{\frac{\sum\cos 2\pi x_i}{D}}+20+e$ \\
Griewank     & M & NS & $[-600,600]^D$ & $\frac{\sum x_i^2}{4000}-\prod\cos\frac{x_i}{\sqrt{i}}+1$ \\
Schwefel 2.26& M & S  & $[-500,500]^D$ & $418{,}98D-\sum x_i\sin\sqrt{|x_i|}$ \\
Lévy         & M & NS & $[-10,10]^D$   & $\sin^2(\pi w_1)\!+\!\sum(w_i\!-\!1)^2[1\!+\!10\sin^2(\pi w_{i+1})]\!+\!(w_D\!-\!1)^2[1\!+\!\sin^2(2\pi w_D)]$ \\
\bottomrule
\multicolumn{5}{l}{\scriptsize $w_i=1+(x_i-1)/4$;\ \ 
  ótimo global $x^*\approx(420{,}97,\ldots)$ para Schwefel 2.26,
  $\mathbf{1}$ para Rosenbrock e Lévy, $\mathbf{0}$ para as demais.}
\end{tabular}
\end{table}

\begin{figure}[H]
\centering
\begin{subfigure}[b]{0.19\textwidth}
  \includegraphics[width=\textwidth]{figs/sphere.pdf}
  \caption{\scriptsize Sphere}
\end{subfigure}
\begin{subfigure}[b]{0.19\textwidth}
  \includegraphics[width=\textwidth]{figs/rosenbrock.pdf}
  \caption{\scriptsize Rosenbrock}
\end{subfigure}
\begin{subfigure}[b]{0.19\textwidth}
  \includegraphics[width=\textwidth]{figs/sum_squares.pdf}
  \caption{\scriptsize Sum Sq.}
\end{subfigure}
\begin{subfigure}[b]{0.19\textwidth}
  \includegraphics[width=\textwidth]{figs/dixon_price.pdf}
  \caption{\scriptsize Dixon-P.}
\end{subfigure}
\begin{subfigure}[b]{0.19\textwidth}
  \includegraphics[width=\textwidth]{figs/zakharov.pdf}
  \caption{\scriptsize Zakharov}
\end{subfigure}\\[4pt]
\begin{subfigure}[b]{0.19\textwidth}
  \includegraphics[width=\textwidth]{figs/rastrigin.pdf}
  \caption{\scriptsize Rastrigin}
\end{subfigure}
\begin{subfigure}[b]{0.19\textwidth}
  \includegraphics[width=\textwidth]{figs/ackley.pdf}
  \caption{\scriptsize Ackley}
\end{subfigure}
\begin{subfigure}[b]{0.19\textwidth}
  \includegraphics[width=\textwidth]{figs/griewank.pdf}
  \caption{\scriptsize Griewank}
\end{subfigure}
\begin{subfigure}[b]{0.19\textwidth}
  \includegraphics[width=\textwidth]{figs/schwefel.pdf}
  \caption{\scriptsize Schwefel}
\end{subfigure}
\begin{subfigure}[b]{0.19\textwidth}
  \includegraphics[width=\textwidth]{figs/levy.pdf}
  \caption{\scriptsize Lévy}
\end{subfigure}
\caption{Superfícies de contorno das dez funções benchmark ($D=2$).
Linha superior: unimodais; linha inferior: multimodais.}
\label{fig:benchmark_grid}
\end{figure}

%─────────────────────────────────────────────────────────────────────────────
\section{Metodologia}
\label{sec:metodologia}

\subsection{Otimização de Hiperpar\^{a}metros com Optuna}

Os hiperparâmetros foram ajustados por função e dimensão usando
Optuna~\cite{akiba2019} com amostrador TPE e 101~trials por
combinação. O critério de avaliação foi o melhor fitness ao final
de 100.000~FEs com semente $s=42$. Os espaços de busca foram:
AG~— $p_m\!\in\![0{,}01,\,0{,}20]$,
$\sigma\!\in\![0{,}05,\,0{,}50]$,
$\alpha\!\in\![0{,}10,\,0{,}50]$;
DE~— $F\!\in\![0{,}4,\,1{,}5]$, $\text{CR}\!\in\![0{,}1,\,1{,}0]$;
PSO~— $w\!\in\![0{,}3,\,1{,}0]$, $c_1,c_2\!\in\![0{,}5,\,2{,}5]$;
EPSO~— $p_c\!\in\![0{,}1,\,0{,}9]$, $\tau\!\in\![0{,}05,\,0{,}50]$.

\subsection{Protocolo Experimental}

Seguindo Marcelino et al.~\cite{marcelino2018}: $D\in\{30,50,100\}$,
$N=100$, 51~execuções independentes por configuração, critério de
parada exclusivo em 100.000~FEs, inicialização uniforme no domínio
de cada função. A análise estatística aplica o Teste~T de Welch
($\alpha=5\%$) como teste principal e o Wilcoxon Rank-Sum como
verificação complementar não-paramétrica.

%─────────────────────────────────────────────────────────────────────────────
\section{Resultados e Discuss\~{a}o}
\label{sec:resultados}

As Tabelas~\ref{tab:unimodais} e~\ref{tab:multimodais} apresentam,
para cada par (função,~$D$), o algoritmo vencedor com sua média e
desvio padrão do melhor fitness ao final das 51~execuções.

\begin{table}[H]
\caption{Melhor algoritmo por (Função, $D$) — Funções Unimodais.
Média $\pm$ Desvio Padrão das 51~execuções.}
\label{tab:unimodais}
\centering
\scriptsize
\begin{tabular}{llll}
\toprule
\textbf{Função} & \textbf{D} & \textbf{Alg.} & $f(\mathbf{x}^*)$ \\
\midrule
\multirow{3}{*}{Sphere}
  & 30  & PSO  & $\mathbf{1{,}60{\times}10^{-89} \pm 1{,}14{\times}10^{-88}}$ \\
  & 50  & PSO  & $\mathbf{9{,}21{\times}10^{-34} \pm 6{,}02{\times}10^{-33}}$ \\
  & 100 & EPSO & $\mathbf{5{,}82{\times}10^{-14} \pm 8{,}11{\times}10^{-14}}$ \\
\midrule
\multirow{3}{*}{Rosenbrock}
  & 30  & --- & $\mathbf{---}$ \\
  & 50  & --- & $\mathbf{---}$ \\
  & 100 & --- & $\mathbf{---}$ \\
\midrule
\multirow{3}{*}{Sum Squares}
  & 30  & --- & $\mathbf{---}$ \\
  & 50  & --- & $\mathbf{---}$ \\
  & 100 & --- & $\mathbf{---}$ \\
\midrule
\multirow{3}{*}{Dixon-Price}
  & 30  & --- & $\mathbf{---}$ \\
  & 50  & --- & $\mathbf{---}$ \\
  & 100 & --- & $\mathbf{---}$ \\
\midrule
\multirow{3}{*}{Zakharov}
  & 30  & --- & $\mathbf{---}$ \\
  & 50  & --- & $\mathbf{---}$ \\
  & 100 & --- & $\mathbf{---}$ \\
\bottomrule
\end{tabular}
\end{table}

\begin{table}[H]
\caption{Melhor algoritmo por (Função, $D$) — Funções Multimodais.
Média $\pm$ Desvio Padrão das 51~execuções.}
\label{tab:multimodais}
\centering
\scriptsize
\begin{tabular}{llll}
\toprule
\textbf{Função} & \textbf{D} & \textbf{Alg.} & $f(\mathbf{x}^*)$ \\
\midrule
\multirow{3}{*}{Rastrigin}
  & 30  & --- & $\mathbf{---}$ \\
  & 50  & --- & $\mathbf{---}$ \\
  & 100 & --- & $\mathbf{---}$ \\
\midrule
\multirow{3}{*}{Ackley}
  & 30  & --- & $\mathbf{---}$ \\
  & 50  & --- & $\mathbf{---}$ \\
  & 100 & --- & $\mathbf{---}$ \\
\midrule
\multirow{3}{*}{Griewank}
  & 30  & --- & $\mathbf{---}$ \\
  & 50  & --- & $\mathbf{---}$ \\
  & 100 & --- & $\mathbf{---}$ \\
\midrule
\multirow{3}{*}{Schwefel 2.26}
  & 30  & --- & $\mathbf{---}$ \\
  & 50  & --- & $\mathbf{---}$ \\
  & 100 & --- & $\mathbf{---}$ \\
\midrule
\multirow{3}{*}{Lévy}
  & 30  & --- & $\mathbf{---}$ \\
  & 50  & --- & $\mathbf{---}$ \\
  & 100 & --- & $\mathbf{---}$ \\
\bottomrule
\end{tabular}
\end{table}

\subsection{Curvas de Converg\^{e}ncia M\'{e}dia}

As curvas são obtidas interpolando as 51~trajetórias em uma grade
comum de FEs e calculando a média ponto a ponto com banda de $\pm1\sigma$.
O eixo~$x$ representa FEs acumulados e o eixo~$y$ usa escala simlog.

% ── Inserir figuras geradas pelo notebook ─────────────────────────────────────
% \includegraphics[width=\textwidth]{figs/conv_sphere.pdf} ...

Nas funções unimodais separáveis (Sphere, Sum~Squares), PSO e EPSO
convergem rapidamente em $D=30$, aproveitando a independência entre
dimensões; a vantagem se mantém em $D=100$ para o EPSO pela
auto-adaptação dos pesos. Em Rosenbrock, o vale estreito não-separável
favorece o PSO pela atualização coordenada de posição, embora todos
os algoritmos apresentem convergência lenta. Para Zakharov, o
crescimento polinomial penaliza AG e DE em alta dimensão, com EPSO
mostrando maior resiliência. Nas funções multimodais, Rastrigin e
Ackley revelam a maior separação entre algoritmos: AG mantém
diversidade por mais gerações, chegando a melhores regiões antes
da estagnação. A natureza deceptiva de Schwefel~2.26 — ótimo global
próximo de $x^*\approx420{,}97$ — prejudica algoritmos que inicializam
e exploram a região central do domínio, onde os gradientes locais
apontam para falsos ótimos.

\subsection{Boxplots}

Os boxplots das 51~execuções para Rosenbrock, Rastrigin, Zakharov e
Schwefel~2.26 complementam as tabelas revelando dispersão e outliers
além da média.

% ── Inserir figuras geradas pelo notebook ─────────────────────────────────────
% \includegraphics[width=\textwidth]{figs/box_rosenbrock.pdf} ...

Algoritmos com caixas estreitas e poucos outliers são mais robustos.
Em Schwefel~2.26, AG apresenta caixas notavelmente estreitas: o ótimo
periférico é frequentemente localizado, enquanto DE, PSO e EPSO mostram
alta variância com muitos outliers extremos, indicando que a maioria
das execuções converge para falsos ótimos. Em Rastrigin, AG e DE têm
dispersão menor que PSO e EPSO, sugerindo que a pressão de diversidade
do torneio e da seleção \emph{greedy} individual ajudam a evitar
ótimos locais simétricos.

\subsection{An\'{a}lise Estat\'{i}stica}

A significância é avaliada pelo Teste~T de Welch ($\alpha=5\%$),
seguindo Marcelino et al.~\cite{marcelino2018}: $H_0\!:\mu_1=\mu_2$
para cada par, por função e dimensão. O Wilcoxon Rank-Sum serve
como verificação complementar não-paramétrica, sendo robusto a
distribuições assimétricas e imune ao cancelamento catastrófico
que afeta o Teste~T em casos de convergência próxima à precisão
numérica de \emph{float64} (e.g., Sphere em $D=30$).

\begin{table}[H]
\caption{Ranking de vitórias ($\alpha=5\%$) sobre
$6\,\text{pares}\times30\,\text{configurações}=180\,\text{testes}$.}
\label{tab:ranking}
\centering
\begin{tabular}{lcccc}
\toprule
\textbf{Teste} & \textbf{AG} & \textbf{DE} & \textbf{PSO} & \textbf{EPSO} \\
\midrule
Teste T de Welch  & 31 & \textbf{40} & 24 & 29 \\
Wilcoxon Rank-Sum & 39 & 34 & \textbf{45} & 39 \\
\bottomrule
\end{tabular}
\end{table}

A divergência concentra-se nas funções onde múltiplos algoritmos
atingem valores próximos de zero (Sphere, Sum~Squares), nas quais
o Wilcoxon detecta diferenças que o Teste~T deixa escapar por
instabilidade numérica. Nas funções multimodais, os dois testes
convergem na maioria dos pares. O Teste~T elege DE como vencedor
global (40~vitórias); o Wilcoxon, mais robusto, coloca PSO em
primeiro (45~vitórias), seguido por AG e EPSO empatados (39 cada).

\subsection{An\'{a}lise de Escalabilidade}

% ── Inserir figuras geradas pelo notebook ─────────────────────────────────────
% \includegraphics[width=\textwidth]{figs/escala_unimodais.pdf} ...
% \includegraphics[width=\textwidth]{figs/escala_multimodais.pdf} ...

Funções não-separáveis (Rosenbrock, Zakharov, Dixon-Price) exibem
degradação mais severa com $D$, pois a interação entre dimensões
amplifica a dificuldade exponencialmente. Entre os algoritmos, EPSO
apresenta a menor taxa de degradação em unimodais não-separáveis
graças à auto-adaptação de pesos; AG degrada mais rapidamente por
não possuir mecanismo de escalonamento de passo. Em multimodais,
DE e PSO mantêm posições relativas estáveis com $D$ crescente,
enquanto AG — apesar de vencer em $D=30$ — perde terreno em $D=100$
em Ackley e Griewank.

\subsection{Velocidade de Converg\^{e}ncia}

A AUC da curva de convergência mede a integral numérica de
$f_\text{best}$ ao longo dos FEs; menor AUC indica acúmulo mais
rápido de boas soluções.

\begin{table}[H]
\caption{AUC das curvas de convergência. Menor AUC = convergência
mais rápida. Melhor por linha em negrito.}
\label{tab:auc}
\centering
\scriptsize
\setlength{\tabcolsep}{3pt}
\begin{tabular}{llccccc}
\toprule
\textbf{Função} & \textbf{D} &
\textbf{AG} & \textbf{DE} & \textbf{PSO} & \textbf{EPSO} & \textbf{Melhor} \\
\midrule
\multirow{3}{*}{Sphere}
  & 30  & $3{,}54{\times}10^{7}$ & $9{,}75{\times}10^{7}$ & \cellcolor{green!20}$\mathbf{8{,}67{\times}10^{6}}$ & $3{,}80{\times}10^{7}$ & PSO \\
  & 50  & $5{,}03{\times}10^{7}$ & $2{,}43{\times}10^{8}$ & \cellcolor{green!20}$\mathbf{2{,}26{\times}10^{7}}$ & $8{,}30{\times}10^{7}$ & PSO \\
  & 100 & $2{,}68{\times}10^{8}$ & $9{,}96{\times}10^{8}$ & \cellcolor{green!20}$\mathbf{1{,}01{\times}10^{8}}$ & $2{,}46{\times}10^{8}$ & PSO \\
\midrule
\multirow{3}{*}{Rosenbrock}
  & 30  & $1{,}99{\times}10^{8}$ & $1{,}27{\times}10^{9}$ & \cellcolor{green!20}$\mathbf{1{,}47{\times}10^{8}}$ & $9{,}96{\times}10^{8}$ & PSO \\
  & 50  & $6{,}43{\times}10^{8}$ & $3{,}89{\times}10^{9}$ & \cellcolor{green!20}$\mathbf{3{,}45{\times}10^{8}}$ & $2{,}23{\times}10^{9}$ & PSO \\
  & 100 & $2{,}24{\times}10^{9}$ & $2{,}06{\times}10^{10}$ & \cellcolor{green!20}$\mathbf{1{,}33{\times}10^{9}}$ & $5{,}37{\times}10^{9}$ & PSO \\
\midrule
\multirow{3}{*}{Sum Sq.}
  & 30  & $5{,}47{\times}10^{6}$ & $1{,}21{\times}10^{7}$ & \cellcolor{green!20}$\mathbf{1{,}24{\times}10^{6}}$ & $5{,}24{\times}10^{6}$ & PSO \\
  & 50  & $1{,}21{\times}10^{7}$ & $5{,}35{\times}10^{7}$ & \cellcolor{green!20}$\mathbf{4{,}89{\times}10^{6}}$ & $2{,}03{\times}10^{7}$ & PSO \\
  & 100 & $8{,}44{\times}10^{7}$ & $4{,}09{\times}10^{8}$ & \cellcolor{green!20}$\mathbf{5{,}96{\times}10^{7}}$ & $1{,}11{\times}10^{8}$ & PSO \\
\midrule
\multirow{3}{*}{Dixon-P.}
  & 30  & $5{,}51{\times}10^{8}$ & $1{,}26{\times}10^{9}$ & \cellcolor{green!20}$\mathbf{1{,}36{\times}10^{8}}$ & $5{,}36{\times}10^{8}$ & PSO \\
  & 50  & $1{,}89{\times}10^{9}$ & $8{,}10{\times}10^{9}$ & \cellcolor{green!20}$\mathbf{6{,}11{\times}10^{8}}$ & $1{,}96{\times}10^{9}$ & PSO \\
  & 100 & $8{,}09{\times}10^{9}$ & $5{,}97{\times}10^{10}$ & \cellcolor{green!20}$\mathbf{4{,}92{\times}10^{9}}$ & $1{,}02{\times}10^{10}$ & PSO \\
\midrule
\multirow{3}{*}{Zakharov}
  & 30  & \cellcolor{green!20}$\mathbf{1{,}94{\times}10^{9}}$ & $3{,}44{\times}10^{9}$ & $1{,}58{\times}10^{10}$ & $4{,}08{\times}10^{9}$ & AG \\
  & 50  & $5{,}27{\times}10^{12}$ & $6{,}92{\times}10^{12}$ & \cellcolor{green!20}$\mathbf{4{,}20{\times}10^{12}}$ & $7{,}28{\times}10^{12}$ & PSO \\
  & 100 & $2{,}55{\times}10^{16}$ & $3{,}33{\times}10^{16}$ & \cellcolor{green!20}$\mathbf{2{,}15{\times}10^{16}}$ & $2{,}21{\times}10^{16}$ & PSO \\
\midrule
\multirow{3}{*}{Rastrigin}
  & 30  & \cellcolor{green!20}$\mathbf{1{,}72{\times}10^{6}}$ & $6{,}64{\times}10^{6}$ & $4{,}55{\times}10^{6}$ & $6{,}47{\times}10^{6}$ & AG \\
  & 50  & \cellcolor{green!20}$\mathbf{3{,}63{\times}10^{6}}$ & $2{,}16{\times}10^{7}$ & $8{,}04{\times}10^{6}$ & $1{,}49{\times}10^{7}$ & AG \\
  & 100 & \cellcolor{green!20}$\mathbf{9{,}55{\times}10^{6}}$ & $5{,}95{\times}10^{7}$ & $2{,}11{\times}10^{7}$ & $3{,}10{\times}10^{7}$ & AG \\
\midrule
\multirow{3}{*}{Ackley}
  & 30  & \cellcolor{green!20}$\mathbf{6{,}04{\times}10^{4}}$ & $1{,}39{\times}10^{5}$ & $1{,}20{\times}10^{5}$ & $2{,}96{\times}10^{5}$ & AG \\
  & 50  & \cellcolor{green!20}$\mathbf{7{,}20{\times}10^{4}}$ & $2{,}12{\times}10^{5}$ & $1{,}91{\times}10^{5}$ & $3{,}44{\times}10^{5}$ & AG \\
  & 100 & $1{,}70{\times}10^{5}$ & $4{,}37{\times}10^{5}$ & $3{,}96{\times}10^{5}$ & \cellcolor{green!20}$\mathbf{1{,}58{\times}10^{5}}$ & EPSO \\
\midrule
\multirow{3}{*}{Griewank}
  & 30  & $3{,}23{\times}10^{5}$ & $1{,}10{\times}10^{6}$ & \cellcolor{green!20}$\mathbf{9{,}93{\times}10^{4}}$ & $5{,}18{\times}10^{5}$ & PSO \\
  & 50  & $9{,}14{\times}10^{5}$ & $2{,}30{\times}10^{6}$ & \cellcolor{green!20}$\mathbf{2{,}57{\times}10^{5}}$ & $9{,}99{\times}10^{5}$ & PSO \\
  & 100 & $4{,}20{\times}10^{6}$ & $8{,}90{\times}10^{6}$ & \cellcolor{green!20}$\mathbf{8{,}59{\times}10^{5}}$ & $2{,}20{\times}10^{6}$ & PSO \\
\midrule
\multirow{3}{*}{Schwefel}
  & 30  & \cellcolor{green!20}$\mathbf{-6{,}51{\times}10^{39}}$ & $1{,}96{\times}10^{8}$ & $3{,}68{\times}10^{8}$ & $4{,}53{\times}10^{8}$ & AG \\
  & 50  & \cellcolor{green!20}$\mathbf{-5{,}10{\times}10^{36}}$ & $6{,}35{\times}10^{8}$ & $7{,}22{\times}10^{8}$ & $9{,}12{\times}10^{8}$ & AG \\
  & 100 & \cellcolor{green!20}$\mathbf{-1{,}29{\times}10^{41}}$ & $2{,}80{\times}10^{9}$ & $1{,}90{\times}10^{9}$ & $2{,}23{\times}10^{9}$ & AG \\
\midrule
\multirow{3}{*}{Lévy}
  & 30  & $1{,}37{\times}10^{5}$ & $3{,}59{\times}10^{5}$ & \cellcolor{green!20}$\mathbf{3{,}20{\times}10^{4}}$ & $2{,}10{\times}10^{5}$ & PSO \\
  & 50  & $1{,}62{\times}10^{5}$ & $1{,}01{\times}10^{6}$ & \cellcolor{green!20}$\mathbf{1{,}36{\times}10^{5}}$ & $4{,}58{\times}10^{5}$ & PSO \\
  & 100 & $6{,}31{\times}10^{5}$ & $4{,}49{\times}10^{6}$ & \cellcolor{green!20}$\mathbf{3{,}78{\times}10^{5}}$ & $1{,}72{\times}10^{6}$ & PSO \\
\midrule
\multicolumn{2}{l}{\textbf{Vitórias AUC} $D=30$}  & 4 & 0 & \textbf{6} & 0 & \\
\multicolumn{2}{l}{\textbf{Vitórias AUC} $D=50$}  & 3 & 0 & \textbf{7} & 0 & \\
\multicolumn{2}{l}{\textbf{Vitórias AUC} $D=100$} & 2 & 0 & \textbf{7} & 1 & \\
\bottomrule
\multicolumn{7}{l}{\scriptsize $^\dagger$AUC negativa (Schwefel/AG):
  fitness negativo por construção da função deslocada.}
\end{tabular}
\end{table}

PSO lidera em velocidade de convergência em 6 a 7 das 10~funções em
todas as dimensões, com vantagem especialmente marcante em funções
separáveis. AG destaca-se em Rastrigin, Ackley e Schwefel, onde a
diversidade genética permite exploração mais efetiva de paisagens
multimodais. DE e EPSO não lideram em nenhuma configuração em termos
de AUC, embora EPSO assuma a liderança em Ackley/$D=100$ — único
caso onde a auto-adaptação de pesos compensa o custo de duas
avaliações por partícula por geração.

%─────────────────────────────────────────────────────────────────────────────
\section{Conclus\~{a}o}
\label{sec:conclusao}

Este trabalho comparou AG, DE, PSO e EPSO em dez funções benchmark
escaláveis ($D\in\{30,50,100\}$) com hiperparâmetros ajustados por
Optuna e 51~execuções independentes por configuração.

Os resultados confirmam a inexistência de um algoritmo dominante: PSO
lidera em funções separáveis e velocidade de convergência (AUC);
AG destaca-se em funções multimodais densas (Rastrigin, Ackley,
Schwefel); EPSO é competitivo em cenários de alta dimensão com
não-separabilidade; DE apresenta robustez estatística mediana pelo
Teste~T mas cede posição ao Wilcoxon. A análise de escalabilidade
mostra que todos os algoritmos degradam com $D$, sendo a degradação
mais severa em funções não-separáveis.

Como limitações, o estudo emprega funções clássicas sem ruído e
população fixa. Trabalhos futuros incluem extensão para DEEPSO e
C-DEEPSO~\cite{marcelino2018}, avaliação em funções CEC com rotação
e deslocamento, e aplicação a problemas reais de sistemas de energia.

%─────────────────────────────────────────────────────────────────────────────
\begin{thebibliography}{9}

\bibitem{goldberg1989}
Goldberg, D.E.: Genetic Algorithms in Search, Optimization and
Machine Learning. Addison-Wesley, Reading, MA (1989)

\bibitem{storn1997}
Storn, R., Price, K.: Differential evolution -- a simple and efficient
heuristic for global optimization over continuous spaces.
Journal of Global Optimization \textbf{11}(4), 341--359 (1997).
\doi{10.1023/A:1008202821328}

\bibitem{kennedy1995}
Kennedy, J., Eberhart, R.: Particle swarm optimization.
In: Proc.\ IEEE ICNN, vol.~4, pp.~1942--1948. IEEE (1995).
\doi{10.1109/ICNN.1995.488968}

\bibitem{miranda2002}
Miranda, V., Fonseca, N.: EPSO -- evolutionary particle swarm
optimization, a new algorithm with applications in power systems.
In: Proc.\ IEEE/PES T\&D, vol.~2, pp.~745--750. IEEE (2002).
\doi{10.1109/TDC.2002.1177567}

\bibitem{marcelino2018}
Marcelino, C.G. et al.: Solving security constrained optimal power
flow problems: a hybrid evolutionary approach.
Applied Intelligence \textbf{48}(10), 3672--3690 (2018).
\doi{10.1007/s10489-018-1167-5}

\bibitem{akiba2019}
Akiba, T. et al.: Optuna: A next-generation hyperparameter
optimization framework. In: Proc.\ KDD, pp.~2623--2631. ACM (2019).
\doi{10.1145/3292500.3330701}

\bibitem{wolpert1997}
Wolpert, D.H., Macready, W.G.: No free lunch theorems for
optimization. IEEE Trans.\ Evol.\ Comput.\ \textbf{1}(1),
67--82 (1997). \doi{10.1109/4235.585893}

\bibitem{molga2005}
Molga, M., Smutnicki, C.: Test functions for optimization needs.
Technical Report (2005).

\end{thebibliography}

\end{document}